\section{Large idempotents with ordered annihilations}
\label{op:modular-section}

The graph has supplied the ratio $t/m<1/100$. Our present objective is
an idempotent in the finite group algebra $K[H]$ whose regular rank
fraction exceeds $1/96$. It must also have two prescribed annihilations,
on opposite sides. Here
\[
 \ell=(p-1)/2,\qquad D=3(p-1),\qquad w=|H|.
\]
The \emph{regular rank} of $e\in K[H]$ means
$\dim_K eK[H]$, the dimension of the image of left multiplication.
Although this map multiplies on the left, its image is a right ideal.

\begin{proposition}\label{op:modular}
There is a specified idempotent $e\in K[H]$ satisfying, for every $s\in P$,
\[
 (u(s)-1)^\ell e=0,\qquad e(v(s)-1)^\ell=0.
\]
Its regular rank $d$ satisfies $d/w>1/96$. Its coefficients are obtained
by finite operations in the specified field $K$.
\end{proposition}

\subsection{The annihilated column space}

On $K^P$, let $\mathbf e_x$ be the standard basis and
$L_s\mathbf e_x=\mathbf e_{s+x}$. For the three standard generators
$\varepsilon_j$ of $P$ set $X_j=L_{\varepsilon_j}-I$. The $X_j$ commute,
and characteristic $p$ gives $X_j^p=0$. The vectors
\[
 X_1^{n_1}X_2^{n_2}X_3^{n_3}\mathbf e_0,
 \qquad 0\le n_1,n_2,n_3<p,
\]
are a basis. Indeed, every translation of $\mathbf e_0$ is a product of
powers of $I+X_j$ applied to $\mathbf e_0$, so the displayed family spans
$K^P$; its cardinality is the dimension $p^3$.

This is the elementary abelian augmentation filtration in the setting
studied by Jennings~\cite{Jennings1941}. Here the exact degree argument
can be given directly. Let $B$ have as columns the displayed basis
vectors with $n_1+n_2+n_3\ge D-\ell+1$, in order. For every $s$, the
matrix $L_s-I$ is a polynomial in the $X_j$ with zero constant term.
Multiplication by its $\ell$th power raises every column's total degree
to at least $D+1$. Such a monomial has an exponent at least $p$, so
vanishes. Applying the same assertion to $-s$ and transposing yields
\begin{equation}\label{op:B-annihilation}
 (L_s-I)^\ell B=0,\qquad B^{\mathsf T}(L_s-I)^\ell=0.
\end{equation}
The transpose identity uses $L_{-s}^{\mathsf T}=L_s$; no scalar
conjugation is present.

The columns are independent. Their number $h$ can be counted by setting
$a_j=p-1-n_j$: the new nonnegative indices have sum at most $\ell-1$.
Their individual upper bounds are automatic because $\ell-1<p$.
Stars and bars gives
\begin{equation}\label{op:h}
 h=\binom{\ell+2}{3},\qquad
 \frac{h}{p^3}=\frac{(p-1)(p+1)(p+3)}{48p^3}>\frac1{48}.
\end{equation}
For the strict inequality, the numerator minus $p^3$ is
$3p^2-p-3>0$ for $p\ge3$.

The same subspace is killed on the two sides in
Equation~\eqref{op:B-annihilation}, but an arbitrary projection onto it
need not give the outgoing conjugate annihilation. We therefore choose
the row factor by a diagonal weighted pairing.

\subsection{Many nonsingular diagonal pairings}

For a finite function group $A_*=(\mathbb Z/M\mathbb Z)^E$, its
characters $\chi:A_*\to K^*$ are all exponent lists of powers of
$\zeta$. Put
\begin{equation}\label{op:fourier}
 p_\chi^{A_*}=|A_*|^{-1}\sum_{a\in A_*}\chi(a)^{-1}a.
\end{equation}
The denominator is invertible. Reindexing the sum gives
$a p_\chi^{A_*}=\chi(a)p_\chi^{A_*}$. A nontrivial character has sum
zero, since translating its sum by an element on which it is nontrivial
multiplies it by a scalar different from one. It follows that
\[
 p_\chi^{A_*}p_\psi^{A_*}=\delta_{\chi\psi}p_\chi^{A_*},\qquad
 \sum_\chi p_\chi^{A_*}=1.
\]
For the last identity, sum coefficients over the independent character
exponents: every nonzero coordinate of a nonidentity element is detected
by a power of $\zeta$. Each idempotent is nonzero, there are $|A_*|$ of
them, and orthogonality makes them independent. They form a basis of
$K[A_*]$.

Specialize to $A_0$. Its character action is
$(s\lambda)(a)=\lambda((-s)\cdot a)$. Define
\[
 \chi_x=x\lambda,\qquad
 D_\lambda=\diag(\chi_x(c):x\in P),\qquad
 N_\lambda=B^{\mathsf T}D_\lambda^{-1}B.
\]
Call $\lambda$ admissible if its stabilizer in $P$ is trivial and
$N_\lambda$ is invertible. This condition is a finite test.

\begin{lemma}\label{op:good-count}
More than half of the $|A_0|=M^{p^3}$ characters are admissible.
\end{lemma}
\begin{proof}
For each nonzero $s\in P$, translation partitions $P$ into $p^2$ cycles
of length $p$. A character fixed by $s$ has constant exponent on each
cycle, so exactly $M^{p^2}$ characters are fixed. A union bound therefore
bounds the fraction with nontrivial stabilizer by
\[
 (p^3-1)M^{p^2-p^3}.
\]

Let $c_r$ be the coordinate generator at $r$. Since
$\chi_x(c)=\lambda(c_{-x})$, the diagonal entries of $D_\lambda^{-1}$
range independently over $K^*$ as $\lambda$ ranges over all characters.
For independent variables $\tau_x$, Cauchy--Binet gives
\begin{equation}\label{op:weighted-gram}
 Q(\tau):=\det\bigl(B^{\mathsf T}\diag(\tau_x)B\bigr)
 =\sum_{\substack{J\subseteq P\\|J|=h}}\det(B_J)^2\prod_{x\in J}\tau_x.
\end{equation}
Here $B_J$ consists of the indicated rows in their fixed order.
Some maximal minor is nonzero because $B$ has full column rank. The
monomials for distinct $J$ are distinct, so $Q$ is nonzero of total
degree $h$, even though the unweighted Gram matrix could be singular.

We use the polynomial grid estimate of
Schwartz~\cite[Lemma~1 and Corollary~1]{Schwartz1980}; see also the
related identity-testing method of Zippel~\cite{Zippel1979}. A nonzero
polynomial of total degree at most $a$ vanishes on at most an $a/M$
fraction of a product grid whose coordinate sets each have size $M$.
To prove this, induct on the number of variables. If the degree in the
last variable is $b$, its leading coefficient has degree at most $a-b$.
It vanishes at at most an $(a-b)/M$ fraction of the earlier choices.
For each remaining choice, the last-variable polynomial has at most
$b$ roots, contributing at most $b/M$. The sum proves the bound; the
constant case starts the induction.

Applying the estimate on $(K^*)^{p^3}$ to
Equation~\eqref{op:weighted-gram}, the singular-pairing fraction is at
most $h/M$. Both excluded fractions are strictly smaller than $1/4$:
\[
 \frac hM\le\frac{p^3}{p^4-1}<\frac14,\qquad
 (p^3-1)M^{-p^2(p-1)}\le\frac{p^3-1}{M}<\frac14.
\]
The estimates use $p>600$; for example,
$p^4-1-4p^3=p^3(p-4)-1>0$. Their union therefore excludes fewer than
half of the characters.
\end{proof}

Select the first admissible character from each $P$-orbit containing one,
and call the set of selected characters $\mathscr R$. Each selected
orbit is free, of size $p^3$. These orbits cover all admissible characters,
so
\begin{equation}\label{op:orbit-count}
 |\mathscr R|>\frac{|A_0|}{2p^3}.
\end{equation}
This argument does not need admissibility to be constant on an orbit.

\subsection{Realizing the projections in the group algebra}

Fix $\lambda\in\mathscr R$. For $x,y\in P$, define
\[
 E_{xy}^{\lambda}=p_{x\lambda}^{A_0}u(x-y).
\]
Conjugation gives $u(s)p_\chi^{A_0}u(s)^{-1}=p_{s\chi}^{A_0}$.
Fourier orthogonality and freeness of the orbit therefore imply
\[
 E_{xy}^{\lambda}E_{zt}^{\lambda}=\delta_{yz}E_{xt}^{\lambda}.
\]
Indeed the Fourier product can survive only when
$x\lambda=(x-y+z)\lambda$, which is equivalent to $y=z$.
Products between different orbits vanish. The sum
$q_\lambda=\sum_x p_{x\lambda}^{A_0}$ is central: it commutes with
$A_0$ and its summands are permuted by $P$. The basis
$p_{x\lambda}^{A_0}u(s)$ of its ideal is indexed equally by $(x,y)$
through $s=x-y$. Thus these units prove
\[
 q_\lambda K[H]\cong\Mat_{p^3}(K).
\]
This is an explicit character-orbit block, in the normal-subgroup
representation setting of Clifford~\cite{Clifford1937Announcement}.
It does not assert semisimplicity of the entire modular algebra.

On this block, the elements $u(s)$, $c$, and $v(s)$ act respectively by
\[
 L_s,\qquad D_\lambda,\qquad D_\lambda L_sD_\lambda^{-1}.
\]
For example, left multiplication by $u(s)$ shifts row $x$ to row $s+x$,
whereas $c$ multiplies row $x$ by $(x\lambda)(c)$.

Define a matrix and then a finite group-algebra element by
\begin{equation}\label{op:explicit-e}
 T_\lambda=BN_\lambda^{-1}B^{\mathsf T}D_\lambda^{-1},\qquad
 e=\sum_{\lambda\in\mathscr R}\sum_{x,y\in P}
       (T_\lambda)_{xy}E_{xy}^{\lambda}.
\end{equation}
The inverse exists by the selection criterion. Since
$B^{\mathsf T}D_\lambda^{-1}B=N_\lambda$, we have $T_\lambda B=B$
and $T_\lambda^2=T_\lambda$. Its image is exactly $\im B$, of dimension
$h$. Orthogonality of the blocks proves $e^2=e$; on the complementary
ideal its component is zero.

We now compute the regular rank, rather than the smaller rank of a
single representation matrix. In $\Mat_{p^3}(K)$, left multiplication
by $T_\lambda$ has image the matrices whose $p^3$ columns all lie in
$\im B$. Its image dimension is $hp^3$. Summing the selected blocks gives
\begin{equation}\label{op:regular-rank}
 d=hp^3|\mathscr R|,\qquad
 \frac d w=\frac{h|\mathscr R|}{|A_0|}
 >\frac{h}{2p^3}>\frac1{96}.
\end{equation}

The incoming annihilation follows from
$(L_s-I)^\ell T_\lambda=0$. For the outgoing annihilation, keep the
factors in their actual order:
\[
 \begin{split}
 B^{\mathsf T}D_\lambda^{-1}
     (D_\lambda L_sD_\lambda^{-1}-I)^\ell
 &=B^{\mathsf T}(L_s-I)^\ell D_\lambda^{-1}\\
 &=0.
 \end{split}
\]
Equation~\eqref{op:B-annihilation} proves both assertions on every
selected block. They also hold on the complementary ideal, where $e$
is zero. This proves Proposition~\ref{op:modular}, including its finite
prescription and the strict rank estimate.

\subsection{From opposite annihilations to an ordered zero product}

Take copies $H_i$ of $H$. For each edge $i<j$, identify
$v(1,\bar j,\bar j^2)$ in $H_i$ with $u(1,\bar i,\bar i^2)$ in $H_j$,
where bars denote residues in $\F_p$. Let $G_0$ be the resulting
free-product quotient, and let $e_i$ denote the image of $e$ from $H_i$.
The injectivity of these copies is proved in Section~\ref{sec:gluing}.
The following ordered identity
already follows from the local calculations and the quotient relations.

\begin{corollary}\label{op:edge-zero}
For every edge $i<j$ of $\Gamma$, $e_i e_j=0$ in $K[G_0]$.
\end{corollary}
\begin{proof}
Let $y$ be its generator, initially of order $p$ on either local side,
and put $X=y-1$. In the local cyclic group algebra
$K[\langle y\rangle]=K[X]/(X^p)$, the annihilator of right multiplication
by $X^\ell$ is precisely the ideal $(X^{p-\ell})$. Choose representatives
of the left cosets $a\langle y\rangle$ in $H_i$; they give a free right
$K[\langle y\rangle]$-basis of $K[H_i]$. Thus the outgoing identity
$e_iX^\ell=0$, applied coefficientwise in this basis before taking
images in the quotient, gives $e_i=qX^{p-\ell}$ for some local $q$.
On the incoming side $X^\ell e_j=0$. Since $p-\ell=\ell+1$,
\[
 e_i e_j=qX^{\ell+1}e_j=qX(X^\ell e_j)=0.
\]
No factor is commuted across $X$, and no reverse ordered product is
asserted.
\end{proof}

We have constructed the large regular rank and the edge vanishings.
To compare them with the graph factorization, the remaining structural
step must put exact rectangular realizations of all those ranks inside
one nonzero ring.
